\section{NVIDIA is a simulation company}

La sección anterior trató la división del trabajo en la cadena industrial; esta se centra en NVIDIA. Xie Chen (谢晨) recuerda que Jensen Huang (黄仁勋) dijo internamente: NVIDIA is a simulation company. El sentido de esta frase no es solo publicitario: replantea a NVIDIA, de una empresa que «vende GPU», a una empresa que «usa el cómputo para crear mundos en los que se puede entrenar». Los videojuegos son simulación al servicio de la experiencia humana; Omniverse es simulación al servicio de la industria y de los gemelos digitales; Isaac Sim es simulación al servicio de los robots y de la IA física.

El three computer problem de NVIDIA también se relaciona con esto: la primera computadora es el centro de datos; la segunda, el dispositivo de IA física en el extremo; la tercera, la computadora de simulación. El centro de datos entrena los grandes modelos; los robots y vehículos en el extremo ejecutan los modelos; la computadora de simulación genera el mundo físico, los datos y los entornos de evaluación. Cuanto más grave es la escasez de datos de robótica, más se parece la tercera computadora a una puerta de entrada estratégica.

\begin{figure}[H]
\centering
\includegraphics[width=0.96\textwidth]{nvidia-simulation-company.png}
\caption{NVIDIA is a simulation company: la capacidad de cómputo, los motores de física, los entornos de simulación y el entrenamiento de robots forman un ciclo cerrado. Diagrama conceptual propio, elaborado a partir de la conversación de 01:15:33--01:21:25.}
\end{figure}

\begin{knowledgebox}{Lectura de la figura: la estrategia de simulación de NVIDIA conecta tres tipos de cómputo}
A la izquierda de la figura está la capacidad de cómputo para entrenamiento e inferencia; en el centro, los motores de física, el renderizado, los activos y las API; a la derecha, el despliegue en robots y en conducción autónoma. La fortaleza de NVIDIA no es una herramienta aislada, sino unir en una infraestructura las GPU, Omniverse, Isaac Sim, PhysX, Warp/Newton y el ecosistema de robótica. Cuanto más se convierte la simulación en la forma de producir datos, más se acerca una empresa de cómputo a ser una empresa de datos.
\end{knowledgebox}

\subsection{Isaac Sim, MuJoCo y la siguiente generación de motores de física}

Esta subsección estructura un poco el ecosistema de simuladores que aparece en la entrevista. Isaac Sim, de NVIDIA, está construido sobre PhysX y el renderizado de Omniverse; sus ventajas son las GPU, el renderizado y el ecosistema industrial. MuJoCo, vinculado a Google/DeepMind, ha sido durante mucho tiempo un motor de física de uso común en la investigación en robótica; sus ventajas son la física y la API, pero el renderizado no es su punto fuerte. Como el RL de extremo a extremo requiere un paralelismo masivo en GPU, la versión de MuJoCo para CPU difícilmente alcanza el caudal de procesamiento necesario; MJX y MuJoCo Warp/Newton representan la evolución hacia la aceleración por GPU, el ecosistema de código abierto y una base física unificada.

Lo que más vale la pena aprender aquí no es qué motor gana, sino que el simulador se convertirá en una capa de la competencia entre ecosistemas. Cuanto más abierto, eficiente y mantenible sea el motor de física subyacente, más rápido podrán desarrollarse las capas superiores: activos, API, generación de datos y plataformas de RL. Si una empresa como Guanglun (光轮) logra aportar activos, parámetros físicos y cadenas de herramientas, se volverá una parte importante de este ecosistema sin necesidad de poseer en exclusiva el motor subyacente.

\begin{knowledgebox}{Términos clave: las cuatro capas de la pila de simulación}
\begin{tabularx}{\textwidth}{p{0.18\textwidth}p{0.36\textwidth}X}
\toprule
Capa & Contenido representativo & Función \\
\midrule
Physics Solver & PhysX, MuJoCo, Newton, etc. & Calcula la mecánica, las colisiones, las restricciones y el movimiento. \\
Rendering & Cadenas de renderizado como Omniverse & Genera las observaciones visuales y las entradas de los sensores. \\
Sim-ready Assets & Activos interactivos, escenas, cuerpos de robot & Proporcionan la materia prima del mundo de entrenamiento. \\
API/Framework & Exportación de datos, metadata, interfaces de entornos de RL & Permiten que el entrenamiento y la evaluación de modelos invoquen la simulación. \\
\bottomrule
\end{tabularx}
\end{knowledgebox}

\subsection{Datos sintéticos y grandes modelos}

Esta subsección pasa de los robots a los grandes modelos. En la entrevista, Xie Chen sostiene que los grandes modelos de lenguaje natural también dependen cada vez más de datos sintéticos, porque los datos de internet tienden a agotarse y el propio modelo puede generar nuevas tareas, nuevas soluciones y nuevas evaluaciones. En cierto sentido, GPT también es un generador de datos sintéticos: produce lenguaje natural sin cesar, que puede usarse para entrenar, destilar o evaluar otros modelos.

Para los grandes modelos corporeizados, en cambio, los datos sintéticos son más complejos: no solo requieren texto, sino también 3D, la perspectiva del robot, múltiples sensores, trayectorias de acción e interacción física. Los datos sintéticos de lenguaje natural pueden apoyarse en la generación y la verificación por modelos; los datos sintéticos corporeizados tienen que apoyarse en mundos simulados y en la calibración con la física real. Por eso EP109 da a los datos sintéticos más peso que la discusión habitual sobre grandes modelos.

\begin{warningbox}{Los datos sintéticos de lenguaje y los corporeizados no tienen la misma dificultad}
Los datos de lenguaje pueden generarse y revisarse principalmente en el espacio simbólico; los datos corporeizados tienen que pasar por las restricciones del mundo físico. Para saber si la trayectoria de un robot es válida no basta con la lógica textual: hay que considerar el contacto, la mecánica, la seguridad, los actuadores y el despliegue real.
\end{warningbox}

\subsection{Resumen de la sección}

La estrategia de simulación de NVIDIA muestra que la infraestructura de IA del futuro no se limita a los clústeres de entrenamiento y a los chips en el extremo, sino que también incluye la computadora de simulación que genera los mundos de entrenamiento. Para la robótica, la simulación es a la vez una herramienta de producción de datos, un entorno de evaluación de modelos y una puerta de entrada al ecosistema industrial.
